\chapter[Równania Hamiltona--Jacobiego]{Równania Hamiltona--Jacobiego: charakterystyki, rozwiązania lepkościowe i wzór Hopfa--Laxa}
\label{chap:pde23}

\section{Zagadnienie Cauchy'ego i geometryczne znaczenie hamiltonianu}
Niech $0<T<\infty$, $d\in\N$, $d\ge1$,
$g\colon\R^d\to\R$, a $H\colon\R^d\times\R^d\to\R$
będzie hamiltonianem. Rozważamy równanie
\begin{equation}\label{eq:pde23-hj}
 \partial_tu(t,x)+H(x,\nabla_xu(t,x))=0,
 \qquad u(0,x)=g(x).
\end{equation}
W pierwszej części zakładamy gładkość danych, potem będziemy
rozpatrywać funkcje ciągłe, niekoniecznie różniczkowalne.
Gradienty $D_xH$ i $D_pH$ traktujemy jako wektory kolumnowe;
drugi z nich jest gradientem względem zmiennej pędowej $p$.
Macierze Jacobiego działają na wektory kolumnowe.

Równanie~\eqref{eq:pde23-hj} występuje w mechanice
hamiltonowskiej, optymalnym sterowaniu i prawach zachowania.
Nawet jeśli $g$ i $H$ są gładkie, rozwiązanie klasyczne
może utracić regularność po skończonym czasie.
Po przecięciu charakterystyk ich klasyczna konstrukcja
może przestać wyznaczać jednowartościowe rozwiązanie.
Potrzebujemy wówczas pojęcia rozwiązania i zasady
jednoznaczności, które nie zakładają różniczkowalności.
Lokalną metodę charakterystyk poznaliśmy w rozdziale
~\ref{chap:pde-characteristics}; poniżej wyprowadzimy
jej postać właściwą równaniu Hamiltona--Jacobiego.

\begin{definicja}[Rozwiązanie klasyczne]
Niech $H$ i $g$ będą ciągłe. Funkcja
$u\in C([0,T]\times\R^d)\cap C^1((0,T)\times\R^d)$
jest rozwiązaniem klasycznym~\eqref{eq:pde23-hj}, jeżeli
spełnia równanie punktowo dla $0<t<T$ oraz $u(0,x)=g(x)$.
W konstrukcji charakterystyk z danych początkowych
będziemy zakładać dodatkową regularność aż do chwili
$t=0$; sama powyższa definicja jej nie wymaga.
\end{definicja}

\section{Charakterystyki układu Hamiltona i utrata gładkości}
\begin{twierdzenie}[Układ charakterystyczny]
\label{thm:pde23-char}
Niech $H\in C^2(\R^d_x\times\R^d_p)$ i niech
$u\in C^2((0,T)\times\R^d)$ spełnia
$u_t+H(x,\nabla_xu)=0$. Dla $(t_0,x_0)\in(0,T)\times\R^d$
rozważmy rozwiązanie układu
\begin{equation}\label{eq:pde23-hamilton}
 \dot X=D_pH(X,P),\qquad \dot P=-D_xH(X,P),
 \qquad X(t_0)=x_0,\quad P(t_0)=\nabla_xu(t_0,x_0).
\end{equation}
Na całym jego przedziale istnienia zawartym w $(0,T)$
zachodzi $P(t)=\nabla_xu(t,X(t))$ i
\begin{equation}\label{eq:pde23-action}
 \frac{\mathrm d}{\mathrm dt}u(t,X(t))
 =P(t)\cdot D_pH(X(t),P(t))-H(X(t),P(t)).
\end{equation}
Ponadto $H(X(t),P(t))$ jest stały.
\end{twierdzenie}
\begin{proof}
Pole wektorowe w układzie~\eqref{eq:pde23-hamilton}
jest klasy $C^1$, więc lokalne istnienie i jednoznaczność
wynikają z twierdzenia Picarda--Lindelöfa.
Różniczkujemy równanie Hamiltona--Jacobiego względem $x_j$:
\[
 \partial_tu_{x_j}+\sum_{k=1}^d H_{p_k}(x,\nabla_xu)u_{x_jx_k}
 +H_{x_j}(x,\nabla_xu)=0.
\]
Równanie dla gradientu przyjmuje postać
\[
 \partial_t\nabla_xu+
 D_x^2u\,D_pH(x,\nabla_xu)=-D_xH(x,\nabla_xu).
\]
Pole $(t,x)\mapsto D_pH(x,\nabla_xu(t,x))$
jest klasy $C^1$. Rozwiązujemy lokalnie
$\dot Y=D_pH(Y,\nabla_xu(t,Y))$, $Y(t_0)=x_0$,
i kładziemy $Q(t)=\nabla_xu(t,Y(t))$.
Reguła łańcuchowa i powyższa równość dają
\[
 \dot Q=\partial_t\nabla_xu(t,Y)+D_x^2u(t,Y)\dot Y
       =-D_xH(Y,Q).
\]
Para $(Y,Q)$ spełnia zatem układ Hamiltona
z tymi samymi danymi co $(X,P)$.
Jednoznaczność daje lokalnie $Y=X$, $Q=P$.
Zbiór chwil, w których $P(t)=\nabla_xu(t,X(t))$,
jest względnie domknięty przez ciągłość i względnie
otwarty przez powtórzenie powyższego argumentu
z dowolnej takiej chwili. Ponieważ zawiera $t_0$,
jest całym rozpatrywanym przedziałem.
Reguła łańcuchowa daje teraz
\[
 \frac{\mathrm d}{\mathrm dt}u(t,X(t))
 =u_t(t,X(t))+\nabla_xu(t,X(t))\cdot\dot X(t)
 =-H(X,P)+P\cdot D_pH(X,P).
\]
Wreszcie
\[
 \frac{\mathrm d}{\mathrm dt}H(X,P)
 =D_xH(X,P)\cdot D_pH(X,P)
 -D_pH(X,P)\cdot D_xH(X,P)=0.
\]
\end{proof}

\begin{stwierdzenie}[Lokalna rekonstrukcja rozwiązania z charakterystyk]
Niech $H\in C^2(\R^{2d})$, $g\in C^2(\R^d)$
i $y_0\in\R^d$. Istnieje otoczenie punktu $(0,y_0)$,
w którym charakterystyki z danymi
$X(0,y)=y$, $P(0,y)=\nabla g(y)$ wyznaczają
rozwiązanie klasy $C^2$ równania Hamiltona--Jacobiego
z warunkiem początkowym $g$ (dla nieujemnych czasów).
Konstrukcja jest ważna lokalnie wszędzie tam,
gdzie odwzorowanie $(t,y)\mapsto(t,X(t,y))$
jest lokalnym dyfeomorfizmem.
\end{stwierdzenie}
\begin{proof}
Twierdzenie o lokalnym przepływie pola klasy $C^1$
daje funkcje $(X,P)$ klasy $C^1$ zmiennych $(t,y)$
w otoczeniu $(0,y_0)$; można tu rozważać także
ujemne czasy. Definiujemy
\[
 Z(t,y)=g(y)+\int_0^t
 \bigl(P(r,y)\cdot\dot X(r,y)-H(X(r,y),P(r,y))\bigr)
 \,\mathrm d\lambda_1(r).
\]
Funkcja podcałkowa jest klasy $C^1$ względem $(r,y)$,
ponieważ $\dot X=D_pH(X,P)$; stąd $Z$ jest klasy $C^1$.
Kluczowa jest tożsamość
\[
 D_yZ(t,y)=P(t,y)^{\mathsf T}D_yX(t,y),
\]
gdzie pochodna funkcji skalarnej jest wierszem.
Dla $j=1,\ldots,d$ różniczkujemy definicję $Z$:
\[
 \begin{split}
 \partial_{y_j}\dot Z
 &=\partial_{y_j}P\cdot\dot X+
 P\cdot\partial_{y_j}\dot X\\
 &\quad-D_xH\cdot\partial_{y_j}X
          -D_pH\cdot\partial_{y_j}P\\
 &=P\cdot\partial_{y_j}\dot X+
          \dot P\cdot\partial_{y_j}X\\
 &=\partial_t(P\cdot\partial_{y_j}X).
 \end{split}
\]
Pochodna mieszana $\partial_t\partial_{y_j}X$
istnieje i jest równa $\partial_{y_j}\dot X$:
równanie całkowe przepływu można różniczkować
względem $y$, ponieważ pole i przepływ są klasy $C^1$.
W chwili $0$ mamy $\partial_{y_j}Z=g_{y_j}$
oraz $P(0,y)\cdot\partial_{y_j}X(0,y)=g_{y_j}$.
Całkowanie daje żądaną tożsamość.

W chwili $0$ zachodzi $D_yX(0,y_0)=I$.
Twierdzenie o funkcji odwrotnej zastosowane do
$(t,y)\mapsto(t,X(t,y))$ daje lokalną odwrotność
$(t,x)\mapsto(t,Y(t,x))$ klasy $C^1$.
Definiujemy $u(t,x)=Z(t,Y(t,x))$, czyli
$u(t,X(t,y))=Z(t,y)$.
Różniczkując względem $y$, otrzymujemy
\[
 (\nabla_xu(t,X(t,y)))^{\mathsf T}D_yX(t,y)
 =P(t,y)^{\mathsf T}D_yX(t,y).
\]
Odwracalność $D_yX$ daje $\nabla_xu(t,X(t,y))=P(t,y)$.
Różniczkując względem $t$ przy ustalonym $y$,
otrzymujemy $u_t+P\cdot\dot X=\dot Z$,
czyli $u_t=-H(X,P)$.
Po podstawieniu $y=Y(t,x)$ daje to równanie.
Obie pierwsze pochodne $u$ są klasy $C^1$,
ponieważ $P,Y,H$ mają taką regularność.
Wobec tego $u$ jest klasy $C^2$.
Warunek początkowy wynika z $X(0,y)=y$ i $Z(0,y)=g(y)$.
Dowód działa też na późniejszej lokalnej gałęzi
odwrotności przepływu, dopóki charakterystyki
istnieją i macierz $D_yX$ jest odwracalna.
\end{proof}

Dla $H=H(p)$ druga część układu daje $P(t)=P(0)$.
Jeśli $g\in C^2(\R^d)$, charakterystyka wychodząca
z $y$ ma zatem postać
\begin{equation}\label{eq:pde23-explicit-char}
 X(t,y)=y+tD_pH(\nabla g(y)),\qquad P(t,y)=\nabla g(y).
\end{equation}
Wartość rozwiązania na odwracalnej gałęzi wynosi
\[
 Z(t,y)=g(y)+t\bigl[
 \nabla g(y)\cdot D_pH(\nabla g(y))-H(\nabla g(y))\bigr].
\]
Warunek lokalnej odwracalności jest kontrolowany przez
\begin{equation}\label{eq:pde23-jacobian}
 D_yX(t,y)=I+tD_{pp}^2H(\nabla g(y))D^2g(y).
\end{equation}
Równość wynika z reguły łańcuchowej.
Zerowanie się wyznacznika oznacza utratę lokalnej
odwracalności w klasie dyfeomorfizmów $C^1$;
nie oznacza samo przez się utraty różnowartościowości.

\begin{przyklad}[Powstanie osobliwości z gładkich danych]
\label{ex:pde23-cos}
Niech $d=1$, $H(p)=p^2/2$ i $g(y)=\cos y$. Wtedy
$X(t,y)=y-t\sin y$, $P(t,y)=-\sin y$,
$Z(t,y)=\cos y+(t/2)\sin^2 y$,
a $\partial_yX=1-t\cos y$.
Dla $0\le t<1$ mamy $\partial_yX\ge1-t>0$.
Ponadto $|X(t,y)-y|\le t$, więc odwzorowanie
$y\mapsto X(t,y)$ jest różnowartościowe i na $\R$.
Jego odwrotność jest klasy $C^1$ i daje globalne
w przestrzeni rozwiązanie klasy $C^2$ dla $t<1$.
Różniczkując $u_x(t,X(t,y))=-\sin y$
względem $y$, dostajemy
\[
 u_{xx}(t,X(t,y))=-\frac{\cos y}{1-t\cos y}.
\]
W szczególności $u_{xx}(t,0)=-1/(1-t)$;
rozwiązania nie można przedłużyć w klasie $C^2$
przez punkt $(1,0)$.

W chwili $t=1$ pochodna $X(1,\cdot)$ zeruje się
w $2\pi\mathbb Z$, ale odwzorowanie pozostaje ściśle
rosnące: całka nieujemnej pochodnej po każdym
przedziale dodatniej długości jest dodatnia.
Dla $t>1$ funkcja $y-t\sin y$ ma ujemną pochodną
w zerze i jest dodatnia w $\pi$. Twierdzenie
o wartości pośredniej daje $y_*\in(0,\pi)$
z $X(t,y_*)=X(t,0)=0$.
Pędy na tych dwóch charakterystykach są różne,
gdyż $P(t,0)=0$ i $P(t,y_*)=-\sin y_*\ne0$.
Nie można zatem zachować wszystkich charakterystyk
wychodzących z danych początkowych jako wykresu
jednego gradientu. Nie jest to jeszcze dowód
nieistnienia dowolnego rozwiązania w klasie $C^1$
z samej definicji. Poniżej skonstruujemy rozwiązanie
lepkościowe, które istnieje również dla $t>1$.
\end{przyklad}

\section{Rozwiązania lepkościowe: definicja bez pochodnych rozwiązania}
\begin{definicja}[Rozwiązanie lepkościowe]
\label{def:pde23-visc}
Niech $Q=(0,T)\times\Omega$, gdzie $\Omega\subset\R^d$
jest obszarem, a $H\colon\Omega\times\R^d\to\R$ jest
ciągły. Funkcja $u\in C(Q)$ jest \emph{podrozwiązaniem
lepkościowym} równania $u_t+H(x,\nabla_xu)=0$, jeśli dla
każdego $(t_0,x_0)\in Q$ i każdej funkcji
$\varphi$ klasy $C^1$ w otoczeniu tego punktu, dla której
$u-\varphi$ ma tam lokalne maksimum, zachodzi
\begin{equation}\label{eq:pde23-sub}
 \varphi_t(t_0,x_0)+H(x_0,\nabla_x\varphi(t_0,x_0))\le0.
\end{equation}
\emph{Nadrozwiązanie lepkościowe} to funkcja ciągła, dla której
w każdym punkcie lokalnego minimum $u-\varphi$ zachodzi
odwrotna nierówność $\ge0$ (dla wszystkich takich funkcji
testowych). \emph{Rozwiązanie lepkościowe}
jest jednocześnie podrozwiązaniem i nadrozwiązaniem.
Warunek początkowy $u(0,x)=g(x)$ oznacza, że
$u$ ma ciągłe przedłużenie do chwili $0$ o tej wartości.
Na płaskim torusie stosujemy tę samą definicję
w lokalnych współrzędnych otrzymanych z ilorazu
$\R^d/(2\pi\mathbb Z)^d$; zmiana takiej reprezentacji
jest lokalnie translacją, więc nie zmienia pochodnych.
\end{definicja}

\begin{uwaga}
Można równoważnie stosować funkcje testowe z $C^1(Q)$:
lokalną funkcję mnożymy przez gładką funkcję odcinającą,
równą $1$ w mniejszym otoczeniu punktu i mającą nośnik
w otoczeniu, na którym funkcja jest określona,
a następnie przedłużamy przez zero.
Nie zmienia to zachowania $u-\varphi$ przy punkcie
ani pochodnych funkcji testowej w tym punkcie.
Dodanie do $\varphi$ stałej pozwala założyć
$u(t_0,x_0)=\varphi(t_0,x_0)$ bez zmiany nierówności.
Ekstrema w definicji nie muszą być ścisłe.
\end{uwaga}

\begin{uwaga}
Słowo „lepkościowe” pochodzi z przybliżenia równań pierwszego
rzędu równaniami $u_t+H(x,\nabla_xu)=\eps\Delta u$,
$\eps\downarrow0$. Definicja nie wymaga istnienia takiego
przybliżenia ani pochodnych $u$. Korzysta wyłącznie
z pochodnych funkcji testowej.
Nie utożsamiamy jej automatycznie z rozwiązaniem słabym
w sensie dystrybucyjnym.
\end{uwaga}

\begin{stwierdzenie}[Zgodność z rozwiązaniem klasycznym]
\label{prop:pde23-smooth-visc}
Jeżeli $u\in C^1(Q)$, to $u$ jest rozwiązaniem
lepkościowym równania $u_t+H(x,\nabla_xu)=0$
wtedy i tylko wtedy, gdy spełnia je punktowo.
\end{stwierdzenie}
\begin{proof}
Jeżeli $u$ spełnia równanie punktowo i $u-\varphi$ ma
lokalne ekstremum, to na mocy warunku koniecznego
ekstremum wewnętrznego pochodne $u-\varphi$ znikają:
$\varphi_t=u_t$ i $\nabla_x\varphi=\nabla_xu$.
Dostajemy właściwą nierówność, a nawet równość.
W drugą stronę w ustalonym punkcie $(t_0,x_0)$
bierzemy $\varphi=u$. Różnica $u-\varphi\equiv0$
ma zarazem maksimum i minimum.
Obie nierówności wymuszają równość.
\end{proof}

\section{Zasada porównania na torusie}
Przez $\T^d=\R^d/(2\pi\mathbb Z)^d$ rozumiemy płaski torus,
a przez $\lambda_d$ na $\T^d$ miarę indukowaną przez miarę
Lebesgue'a na $[0,2\pi)^d$, bez normalizacji do jedynki;
w szczególności $\lambda_d(\T^d)=(2\pi)^d$.
Najważniejszym argumentem zapewniającym jednoznaczność
jest \emph{podwajanie zmiennych}. Podajemy pełny dowód
modelowego twierdzenia, w którym brak brzegu pozwala
oddzielić istotę argumentu od technicznych barier.

\begin{twierdzenie}[Porównanie rozwiązań lepkościowych]
\label{thm:pde23-comparison}
Niech $T>0$, a $H\colon\T^d\times\R^d\to\R$ będzie
ciągły i niech istnieje $L>0$ takie, że dla wszystkich
$x,y\in\T^d$, $p\in\R^d$ i dla dostatecznie bliskich
$x,y$, przy użyciu tej samej lokalnej reprezentacji
współrzędnych,
\begin{equation}\label{eq:pde23-HxLip}
 |H(x,p)-H(y,p)|\le L|x-y|(1+|p|).
\end{equation}
Niech $u,v\in C([0,T]\times\T^d)$ będą odpowiednio
podrozwiązaniem i nadrozwiązaniem lepkościowym
$w_t+H(x,\nabla w)=0$ w $(0,T)\times\T^d$.
Jeśli $u(0,x)\le v(0,x)$ dla każdego $x$, to
$u(t,x)\le v(t,x)$ dla wszystkich $(t,x)\in[0,T]\times\T^d$.
\end{twierdzenie}
\begin{proof}
Przypuśćmy, że nierówność jest fałszywa. Wybieramy
$0<\tau<T$ oraz $(t_*,x_*)\in(0,\tau)\times\T^d$
takie, że $u(t_*,x_*)-v(t_*,x_*)>0$.
Wybieramy małe $\delta,\kappa>0$ spełniające
$u(t_*,x_*)-v(t_*,x_*)-\delta t_*-\kappa/(\tau-t_*)>0$.
Dla $\eps,\alpha>0$ maksymalizujemy dla
$t\in[0,\tau)$, $s\in[0,\tau]$ i $x,y\in\T^d$ funkcję
\begin{equation}\label{eq:pde23-doubling}
 \Phi_{\eps,\alpha}(t,x,s,y)=u(t,x)-v(s,y)
 -\frac{d_{\T}(x,y)^2}{2\eps}
 -\frac{(t-s)^2}{2\alpha}
 -\delta t-\frac{\kappa}{\tau-t}.
\end{equation}
Z powodu ostatniego składnika maksimum jest osiągane
przy $t<\tau$; przy ustalonym $t$ wybór $s\in[0,\tau]$
nie wymaga osobnej bariery. Maksimum jest dodatnie,
bo można wstawić $t=s=t_*$ i $x=y=x_*$.
Niech $C_0=2\max(\|u\|_\infty,\|v\|_\infty)+1$.
W punkcie dodatniego maksimum wszystkie składniki kary
są nieujemne, toteż $\kappa/(\tau-t_\eps)\le C_0$.
A zatem $t_\eps\le\tau-\kappa/C_0$ niezależnie od
$\eps,\alpha$. Po wykazaniu, że
$|t_\eps-s_\alpha|\to0$, dostajemy również $s_\alpha<\tau$
dla dostatecznie małych parametrów. Odrębnie poniżej
wykluczamy ich zerowanie się przy $t=0$.

Niech $(t_\eps,x_\eps,s_\alpha,y_\eps)$ będzie
punktem maksimum; dla czytelności pomijamy w indeksach
część parametrów. Porównanie z punktem
$(t_\eps,x_\eps,s_\alpha,x_\eps)$ oraz jednostajna
ciągłość $v$ dają
$\frac{d_{\T}(x_\eps,y_\eps)^2}{2\eps}
\le \omega_v(d_{\T}(x_\eps,y_\eps))$.
Z ograniczoności $u,v$ wynika
$d_{\T}(x_\eps,y_\eps)\to0$, a z powyższej
nierówności
\begin{equation}\label{eq:pde23-pen-small}
 \frac{d_{\T}(x_\eps,y_\eps)^2}{\eps}\longrightarrow0.
\end{equation}
Analogicznie porównanie z $s=t$ i jednostajna ciągłość
$v$ w czasie dają $|t_\eps-s_\alpha|\to0$ i
$(t_\eps-s_\alpha)^2/\alpha\to0$. Wybieramy
kolejność granic $\eps\downarrow0$, $\alpha\downarrow0$
lub wspólną kolejność zbieżną do zera.

Dla dostatecznie małych parametrów zarówno $t_\eps$,
jak i $s_\alpha$ należą do $(0,\tau)$.
Gdyby na przykład $t_\eps=0$ dla pewnego ciągu parametrów,
to $s_\alpha\to0$, $d_{\T}(x_\eps,y_\eps)\to0$;
jednostajna ciągłość i warunek początkowy dawałyby
graniczną wartość maksimum $\le0$, w sprzeczności
z ustaloną dodatnią wartością w punkcie $(t_*,x_*)$.
Tak samo wykluczamy $s_\alpha=0$.

Ponieważ punkty przestrzenne są bliskie, kwadrat
odległości ma w ich otoczeniu gładką postać euklidesową.
Oznaczmy $p=(x_\eps-y_\eps)/\eps$ w lokalnych
współrzędnych. Z definicji maksimum skonstruowane
przez zamrożenie $(s,y)$ funkcje testowe dotykają $u$
z góry, a przez zamrożenie $(t,x)$ dotykają $v$ z dołu.
Otrzymujemy zatem
\begin{align*}
 \frac{t_\eps-s_\alpha}{\alpha}+\delta
  +\frac{\kappa}{(\tau-t_\eps)^2}+H(x_\eps,p)&\le0,\\
 \frac{t_\eps-s_\alpha}{\alpha}+H(y_\eps,p)&\ge0.
\end{align*}
Po odjęciu dostajemy
$\delta\le H(y_\eps,p)-H(x_\eps,p)$.
Założenie~\eqref{eq:pde23-HxLip} pozwala ograniczyć
prawą stronę przez
$L|x_\eps-y_\eps|+L|x_\eps-y_\eps|^2/\eps$,
co dąży do zera na mocy~\eqref{eq:pde23-pen-small}.
Sprzeczność dowodzi twierdzenia dla $t<\tau$,
a dowolność $\tau$ i ciągłość dają tezę do czasu $T$.
\end{proof}

\begin{wniosek}[Jednoznaczność i kontrakcja]
\label{cor:pde23-contract}
Przy założeniach Twierdzenia~\ref{thm:pde23-comparison}
rozwiązanie lepkościowe z danymi $g\in C(\T^d)$ jest
co najwyżej jedno. Jeżeli $u,v$ są dwoma rozwiązaniami
dla danych $g,h\in C(\T^d)$, to
$\|u(t)-v(t)\|_{L^\infty(\T^d)}\le
\|g-h\|_{L^\infty(\T^d)}$ dla $0\le t\le T$.
\end{wniosek}
\begin{proof}
Równanie nie zależy od wartości $u$, zatem $v+c$
jest rozwiązaniem, ilekroć $v$ nim jest. Stosujemy
porównanie do pary $(u,v+c)$, gdzie
$c=\|g-h\|_\infty$, a następnie zamieniamy rolami $u,v$.
\end{proof}

\section{Transformata Legendre'a i wzór Hopfa--Laxa}
\begin{definicja}[Funkcja sprzężona Legendre'a]
Niech $H\colon\R^d\to\R$ będzie funkcją wypukłą.
Jej funkcję sprzężoną definiujemy wzorem
$L(v)=\sup_{p\in\R^d}(p\cdot v-H(p))$.
Funkcję $H$ nazywamy \emph{superliniową}, gdy
$H(p)/|p|\to\infty$ przy $|p|\to\infty$.
\end{definicja}

\begin{lemat}[Dualność i kontrola prędkości minimalizatorów]
\label{lem:pde23-legendre}
Niech $H\in C^1(\R^d)$ będzie wypukła i superliniowa.
Wtedy $L$ jest skończona, ciągła, wypukła i superliniowa,
a dla każdego $p\in\R^d$ zachodzi
\begin{equation}\label{eq:pde23-duality}
 H(p)=\sup_{v\in\R^d}(p\cdot v-L(v)).
\end{equation}
Jeżeli $g$ jest globalnie $K$-lipschitzowska, to
$y\mapsto g(y)+tL((x-y)/t)$ osiąga minimum, a każda
jego minimalizująca wartość $v=(x-y)/t$ spełnia
$L(v)-K|v|\le L(0)$. W szczególności prędkości
minimalizatorów są ograniczone liczbą zależną tylko od
$H$ i $K$.
\end{lemat}
\begin{proof}
Ponieważ $H$ jest superliniowa, funkcja
$p\mapsto p\cdot v-H(p)$ dąży do $-\infty$ dla
$|p|\to\infty$, więc supremum definiujące $L(v)$
jest skończone i osiągane. Jako supremum funkcji
afinicznych $L$ jest wypukła i dolnie półciągła,
a skończona funkcja wypukła na całej $\R^d$ jest ciągła.
Dla dowolnego $R>0$ i $v\ne0$ wybieramy
$p=R v/|v|$, zatem $L(v)\ge R|v|-\max_{|p|=R}H(p)$.
Po podzieleniu przez $|v|$ i przejściu $|v|\to\infty$
wnioskujemy o superliniowości $L$, ponieważ $R$ jest
dowolne.

Z definicji $L$ mamy $p\cdot v-L(v)\le H(p)$.
Jeśli jednak $v=D H(p)$, wypukłość daje dla każdego $q$
$H(q)\ge H(p)+v\cdot(q-p)$, stąd
$L(v)=p\cdot v-H(p)$; powstaje równość w~\eqref{eq:pde23-duality}.

Dla ustalonych $x,t>0$ i $y=x-tv$ warunek Lipschitza
daje $g(x-tv)\ge g(x)-Kt|v|$. Ponieważ $L$ jest
superliniowa, funkcja $g(x-tv)+tL(v)$ dąży do
$+\infty$ przy $|v|\to\infty$; ciągłość zapewnia
osiągnięcie minimum. Porównując je z próbą $v=0$,
dostajemy
$g(x-tv)+tL(v)\le g(x)+tL(0)$.
Wykorzystanie nierówności Lipschitza daje ostatnią tezę.
\end{proof}

\begin{twierdzenie}[Wzór Hopfa--Laxa]
\label{thm:pde23-hopflax}
Niech $H\in C^1(\R^d)$ będzie wypukła i superliniowa,
a $g\colon\R^d\to\R$ niech będzie funkcją ograniczoną
i globalnie $K$-lipschitzowską. Dla $t>0$ zdefiniujmy
\begin{equation}\label{eq:pde23-hopflax}
 u(t,x)=\min_{y\in\R^d}\left\{
 g(y)+tL\!\left(\frac{x-y}{t}\right)\right\},
 \qquad u(0,x)=g(x).
\end{equation}
Wtedy $u$ jest ograniczona i globalnie Lipschitzowska
na $[0,T]\times\R^d$ dla każdego $T>0$ oraz jest
rozwiązaniem lepkościowym równania
$u_t+H(\nabla u)=0$. Ponadto
\begin{equation}\label{eq:pde23-dpp}
 u(t,x)=\min_{z\in\R^d}
 \left\{u(s,z)+(t-s)L\!\left(\frac{x-z}{t-s}\right)\right\}
 \quad (0\le s<t).
\end{equation}
\end{twierdzenie}
\begin{proof}
\textbf{Krok 1: ograniczenie minimalizatorów i ciągłość.}
Lemat~\ref{lem:pde23-legendre} daje osiąganie minimum
oraz wspólne ograniczenie $|x-y|/t\le R$.
Wybór dowolnego $v$ w równoważnej reprezentacji
$u(t,x)=\min_v\{g(x-tv)+tL(v)\}$ daje
$u(t,x)\le g(x)+tL(0)$, a nierówność Lipschitza
$g(x-tv)\ge g(x)-Kt|v|$ daje
$u(t,x)\ge g(x)+t\min_v(L(v)-K|v|)$.
Wynika stąd $\|u(t)-g\|_\infty\le Ct$.
Przy przesunięciu $x\mapsto x+h$ każda próba $v$
zmienia składnik $g(x-tv)$ o co najwyżej $K|h|$;
po wzięciu minimów otrzymujemy
$|u(t,x+h)-u(t,x)|\le K|h|$.

\textbf{Krok 2: zasada programowania dynamicznego.}
Oznaczmy $a=t-s$, $b=s$ i użyjmy wypukłości $L$.
Dla dowolnych $y,z\in\R^d$ mamy
\[
 aL\!\left(\frac{x-z}{a}\right)
 +bL\!\left(\frac{z-y}{b}\right)
 \ge (a+b)L\!\left(\frac{x-y}{a+b}\right).
\]
Minimalizacja kolejno po $y,z$ daje kierunek „$\ge$”
w~\eqref{eq:pde23-dpp}. Kierunek przeciwny otrzymujemy
przez wybór $z=y+\frac{s}{t}(x-y)$ na prostym odcinku
łączącym $y$ z $x$, kiedy obie prędkości są równe
$(x-y)/t$ i nierówność wypukłości staje się równością.
Przypadek $s=0$ wynika bezpośrednio z definicji.

Stosując wzór dynamiczny na przedziale długości $h>0$
i argument kroku pierwszego do funkcji $u(s,\cdot)$,
dostajemy $|u(s+h,x)-u(s,x)|\le Ch$ ze stałą
zależną tylko od $K,H$. Razem z oszacowaniem
przestrzennym dowodzi to globalnej ciągłości Lipschitza.

\textbf{Krok 3: nierówność dla funkcji dotykającej z góry.}
Niech $u-\varphi$ ma lokalne maksimum równe zero
w $(t,x)$, $t>0$. Ustalmy dowolne $v\in\R^d$.
Z~\eqref{eq:pde23-dpp} dla małych $h>0$ dostajemy
$u(t,x)\le u(t-h,x-hv)+hL(v)$.
Ponieważ $u(t-h,x-hv)\le\varphi(t-h,x-hv)$,
mamy $\varphi(t,x)\le\varphi(t-h,x-hv)+hL(v)$.
Po podzieleniu przez $h$ i przejściu do granicy
$\varphi_t+v\cdot\nabla\varphi\le L(v)$.
Warunek zachodzi dla każdego $v$; dzięki
\eqref{eq:pde23-duality} otrzymujemy
$\varphi_t+H(\nabla\varphi)\le0$.

\textbf{Krok 4: nierówność dla funkcji dotykającej z dołu.}
Niech $u-\varphi$ ma lokalne minimum równe zero
w $(t,x)$. Wybieramy minimalizator $z_h=x-hv_h$
w~\eqref{eq:pde23-dpp}. Wspólne oszacowanie Lipschitza
funkcji $u(t-h,\cdot)$ i Lemat~\ref{lem:pde23-legendre}
dają ograniczenie $|v_h|\le R$.
Z minimalności i dotykania z dołu wynika
$\varphi(t,x)\ge\varphi(t-h,x-hv_h)+hL(v_h)$.
Taylor z resztą $o(h)$ jednostajną dla $|v_h|\le R$
daje
$\varphi_t+v_h\cdot\nabla\varphi\ge L(v_h)+o(1)$.
Ale $L(v_h)-v_h\cdot\nabla\varphi\ge-H(\nabla\varphi)$,
więc $\varphi_t+H(\nabla\varphi)\ge o(1)$.
Przechodzimy do $h\downarrow0$.
Obie nierówności dowodzą rozwiązania lepkościowego.
\end{proof}

\begin{uwaga}[O jednoznaczności na całej przestrzeni]
Dla autonomicznego $H=H(p)$ powyższe rozwiązanie jest
jedynym ograniczonym, globalnie Lipschitzowskim
rozwiązaniem lepkościowym o danych $g$.
Dowód porównania na $\R^d$ otrzymuje się z dowodu
Twierdzenia~\ref{thm:pde23-comparison}, dodając do
funkcji podwajanej zmiennej składnik
$-\beta\sqrt{1+|x|^2}$, który zapewnia osiąganie
maksimum. Dla funkcji Lipschitzowskich gradienty funkcji
testowych w punktach styczności są jednostajnie
ograniczone. Nierówność dla $H$ wynika więc z jego
lokalnej lipschitzowskości w zmiennej pędowej, a dodatkowy
błąd ma rozmiar $O(\beta)$. Najpierw przechodzimy
$\eps,\alpha\downarrow0$, później $\beta\downarrow0$.
Jest to ta sama zasada porównania, z dodatkową lokalizacją
maksimum; nie korzysta z wypukłości $H$.
\end{uwaga}


\begin{wniosek}[Skończony obszar zależności we wzorze Hopfa--Laxa]
\label{cor:pde23-hopflax-local}
Niech $H$ spełnia założenia Twierdzenia~\ref{thm:pde23-hopflax},
a $g,h\colon\R^d\to\R$ będą ograniczone i globalnie
$K$-lipschitzowskie. Niech $u_g,u_h$ oznaczają funkcje
zdefiniowane przez wzór Hopfa--Laxa odpowiednio dla $g$ i $h$.
Istnieje $R_{H,K}<\infty$, zależne tylko od $H,K$, takie, że
jeśli $g=h$ na domkniętej kuli $\overline B_R(x_0)$, to
\[
 u_g(t,x)=u_h(t,x)
 \quad\text{dla }t>0,\qquad |x-x_0|+tR_{H,K}\le R.
\]
Jeżeli $R_{H,K}=0$, nierówność interpretujemy bez
ograniczenia czasu. Analogiczna własność zachodzi w chwili $t=0$.
\end{wniosek}
\begin{proof}
Z Lematu~\ref{lem:pde23-legendre} wszyscy minimalizatorzy
dla obu danych są postaci $y=x-tv$, gdzie
$L(v)-K|v|\le L(0)$. Ponieważ $L$ jest superliniowa,
ten podpoziom jest zwarty; wybieramy $R_{H,K}$ nie mniejsze
niż maksimum $|v|$ na nim. Gdy
$|x-x_0|+tR_{H,K}\le R$, każdy minimalizator dla $g$
i każdy minimalizator dla $h$ leży w $\overline B_R(x_0)$.
Na tej kuli obie funkcje minimalizowane są identyczne.
Minimum globalne każdej jest osiągane wewnątrz kuli,
więc obie wartości minimalne są równe.
\end{proof}


\section{Rozwiązanie na przeciętych charakterystykach i półwklęsłość}
\begin{twierdzenie}[Półwklęsłość w przypadku kwadratowym]
\label{thm:pde23-semiconcave}
Niech $H(p)=|p|^2/2$, $g$ będzie ograniczona i
Lipschitzowska, a $u$ będzie dane przez
\eqref{eq:pde23-hopflax}. Dla każdego $t>0$
funkcja $x\mapsto u(t,x)-|x|^2/(2t)$ jest wklęsła.
W szczególności $u(t,\cdot)$ jest lokalnie Lipschitzowska
i dwukrotnie różniczkowalna prawie wszędzie w sensie
Alexandrowa (ostatni fakt jest klasycznym twierdzeniem
o funkcjach półwklęsłych, którego tutaj nie dowodzimy).
\end{twierdzenie}
\begin{proof}
Ponieważ $L(v)=|v|^2/2$, mamy
\[
 u(t,x)-\frac{|x|^2}{2t}
 =\inf_{y\in\R^d}\left\{
 g(y)+\frac{|y|^2}{2t}-\frac{x\cdot y}{t}\right\}.
\]
Dla ustalonego $y$ prawa strona pod nawiasem jest
funkcją afiniczną względem $x$. Infimum funkcji
afinicznych jest funkcją wklęsłą. Lokalna
lipschitzowskość wynika tu również bezpośrednio
z Twierdzenia~\ref{thm:pde23-hopflax}.
\end{proof}

\begin{przyklad}[Utrata różniczkowalności i wybór lepkościowy]
\label{ex:pde23-shock}
Dla $g(x)=\cos x$ i $H(p)=p^2/2$ otrzymujemy
$u(t,x)=\min_y\{\cos y+(x-y)^2/(2t)\}$.
Jeśli $0<t<1$, mapa charakterystyk
$y\mapsto y-t\sin y$ jest ściśle rosnąca.
Gdy $t>1$ i $x=0$, funkcja minimalizowana jest parzysta;
punkt $y=0$ jest lokalnym maksimum, ponieważ druga pochodna
w $0$ wynosi $-1+1/t<0$. Dla $t$ dostatecznie bliskiego $1$
ma ona dwa symetryczne globalne minima $\pm a(t)$,
$a(t)>0$, spełniające $a(t)=t\sin a(t)$.
Dla $x$ dodatnich bliskich zera minimalizatorem jest
gałąź $+a(t)$, a dla ujemnych gałąź $-a(t)$.
Z twierdzenia o różniczkowaniu wartości minimalnej
(lub bezpośrednio z różnic ilorazowych) pochodne jednostronne
$u_x(t,0+)$ i $u_x(t,0-)$ wynoszą odpowiednio
$-a(t)/t$ i $a(t)/t$. Funkcja jest więc ciągła, lecz
nie jest różniczkowalna w $x=0$. Rozwiązanie lepkościowe
istnieje mimo przecięcia charakterystyk i jest jednoznaczne
w klasie ograniczonych rozwiązań Lipschitzowskich.
\end{przyklad}

\section{Związek z równaniem Burgersa i warunkiem entropii}
Jeżeli $u$ jest gładka i spełnia
$u_t+(u_x)^2/2=0$, to $w=u_x$ spełnia
$w_t+(w^2/2)_x=0$. Po powstaniu osobliwości
różniczkowanie można interpretować dystrybucyjnie.
Pochodna rozwiązania lepkościowego odpowiada
rozwiązaniu entropijnemu równania Burgersa dla
odpowiednich danych Lipschitzowskich.
Warunek entropii jest więc odpowiednikiem wyboru
lepkościowego na poziomie prawa zachowania.
Nie wyprowadzamy tutaj pełnego twierdzenia Krużkowa;
zostanie ono udowodnione w rozdziale o prawach zachowania.

\begin{stwierdzenie}[Różniczkowanie lepkiego równania Hamiltona--Jacobiego]
\label{prop:pde-v43-hj-viscous}
Niech $f\in C^2(\R)$, $\varepsilon>0$ oraz
$U\in C^{1,3}((0,T)\times\R;\R)$ spełnia
$U_t+f(U_x)=\varepsilon U_{xx}$ punktowo.
Wówczas $w=U_x\in C^{1,2}$ spełnia
$w_t+\partial_x f(w)=\varepsilon w_{xx}$.
Ponadto dla każdej wypukłej $\eta\in C^2(\R)$ i
strumienia $q$ określonego przez $q'=\eta' f'$ zachodzi
\[
 \partial_t\eta(w)+\partial_x q(w)
 =\varepsilon\partial_{xx}\eta(w)
 -\varepsilon\eta''(w)|w_x|^2
 \le\varepsilon\partial_{xx}\eta(w).
\]
\end{stwierdzenie}
\begin{proof}
Różniczkujemy równanie dla $U$ względem $x$:
$w_t+f'(w)w_x=\varepsilon w_{xx}$.
Po pomnożeniu przez $\eta'(w)$ otrzymujemy
$\partial_t\eta(w)+\partial_xq(w)
=\varepsilon\eta'(w)w_{xx}$.
Reguła łańcuchowa daje
$\eta'(w)w_{xx}=\partial_{xx}\eta(w)-\eta''(w)|w_x|^2$.
Ostatni wyraz jest niedodatni przez wypukłość $\eta$.
\end{proof}

\begin{uwaga}[Różniczkowanie nie jest jeszcze przejściem do granicy]
Powyższa tożsamość obowiązuje dla gładkich rozwiązań
\emph{z tą samą lepkością} $\varepsilon$ i łączy
rachunek w rozdziałach~\ref{chap:pde23} i~\ref{chap:pde24}.
Aby stwierdzić, że $\partial_xU$ jest rozwiązaniem
entropijnym dla granicznego $U$, trzeba ponadto
uzasadnić zbieżność pochodnych przy
$\varepsilon\downarrow0$ i zgodność danych początkowych;
sama zbieżność jednostajna $U^\varepsilon\to U$
nie zapewnia silnej zbieżności $U_x^\varepsilon$.
\end{uwaga}


\begin{uwaga}[Zakres wyników]
Twierdzenie o porównaniu udowodniliśmy na torusie
w założeniach~\eqref{eq:pde23-HxLip}; wzór Hopfa--Laxa
udowodniliśmy dla autonomicznego hamiltonianu wypukłego
i superliniowego. Ogólna teoria dla hamiltonianów
niewypukłych, zależnych od $x$ i sterowania optymalnego
wymaga osobnych hipotez i metod. Samo stosowanie
wzoru~\eqref{eq:pde23-hopflax} poza tymi założeniami
nie jest uzasadnione.
\end{uwaga}

\section{Wskazówki do dalszej lektury rozdziału}
W następnym rozdziale przejdziemy od potencjału $u$
do skalarnego prawa zachowania dla $w=u_x$, gdzie
pojawią się rozwiązania słabe, fale uderzeniowe,
warunki Rankine'a--Hugoniota i selekcja entropijna.
Charakterystyki i geometria ich przecinania
pozostaną podstawowym źródłem przykładów.

